\documentclass[12pt]{article}
\usepackage[utf8]{inputenc}
\usepackage[T1]{fontenc}
\usepackage[margin=1in]{geometry}
\usepackage{setspace}
\usepackage{xcolor}
\usepackage[colorlinks=true, urlcolor=blue!60!black, linkcolor=blue!60!black]{hyperref}
\setlength{\parindent}{2em}
\setlength{\parskip}{1.25em}
\renewcommand{\baselinestretch}{1.0}
\begin{document}
\begin{center}
{\LARGE \textit{Luther v. Borden}, 48 U.S. (7 How.) 1 (1849)\par}
\vspace{0.8em}
{\large Case Brief\par}
\end{center}
\vspace{1em}

\section{Memory Jogger}

Federal courts lack authority to decide which of two rival state governments is legitimate under the Guarantee Clause because recognizing a state's lawful government is a nonjusticiable political question committed to Congress and the President.

\section{Detailed Case Facts}

Unlike the other twelve original states, Rhode Island did not adopt a new written constitution after the American Revolution; instead, it continued operating under the royal colonial charter granted by King Charles II in 1663, modified only by legislative acts (\href{https://supreme.justia.com/cases/federal/us/48/1/}{48 U.S. at 35}; \href{https://www.courtlistener.com/opinion/86466/luther-v-borden/}{CourtListener}). The 1663 Charter contained no amendment mechanism and permitted the General Assembly to set voter qualifications, which the legislature restricted to freeholders (landowners), excluding a growing majority of urban and industrial adult male residents (\href{https://supreme.justia.com/cases/federal/us/48/1/}{48 U.S. at 35, 48}).

After repeated petitions to the General Assembly for suffrage reform failed, reform advocates organized Suffrage Associations in 1840--1841, held mass conventions in Newport and Providence, and convened an unauthorized ``People's Convention'' in late 1841 (\href{https://supreme.justia.com/cases/federal/us/48/1/}{48 U.S. at 11--14, 35--36}). That convention drafted the ``People's Constitution,'' extending suffrage to adult white male citizens meeting residency requirements, and conducted an extralegal ratification vote in December 1841 (\href{https://supreme.justia.com/cases/federal/us/48/1/}{48 U.S. at 14, 36}). In January 1842, the convention declared the People's Constitution ratified by a majority of adult male citizens and even a majority of existing freehold voters, and in April 1842 its supporters elected Thomas Wilson Dorr as governor alongside a rival legislature, which briefly convened in Providence on May 3, 1842 (\href{https://supreme.justia.com/cases/federal/us/48/1/}{48 U.S. at 15--17, 36}).

The established charter government, led by Governor Samuel Ward King, repudiated the People's Constitution as an unlawful usurpation and refused to surrender state offices or property (\href{https://supreme.justia.com/cases/federal/us/48/1/}{48 U.S. at 7, 36--37}). When Dorr attempted to seize the state arsenal in Providence in May 1842 and later assembled armed supporters at Chepachet in June 1842, the charter General Assembly enacted ``An Act establishing Martial Law in this State'' on June 25, 1842, and Governor King called out the state militia and applied to President John Tyler for federal assistance (\href{https://supreme.justia.com/cases/federal/us/48/1/}{48 U.S. at 8--9, 37, 44}). Although President Tyler never deployed federal troops, he formally recognized the charter government as Rhode Island's lawful authority and signaled his readiness to call out the militia if necessary---a stance that helped cause Dorr's forces to disperse (\href{https://supreme.justia.com/cases/federal/us/48/1/}{48 U.S. at 44}).

On June 29, 1842, while martial law remained in effect, Luther M. Borden and other members of the fourth regiment of the Rhode Island militia, acting under written orders from Quartermaster John T. Child to arrest insurgents, broke and entered the home of Martin Luther in Warren, Rhode Island, to arrest him for aiding Dorr's movement (\href{https://supreme.justia.com/cases/federal/us/48/1/}{48 U.S. at 2, 10, 34}). Meanwhile, the charter government called its own constitutional convention, whose constitution was ratified and took effect in May 1843, permanently superseding the 1663 Charter (\href{https://supreme.justia.com/cases/federal/us/48/1/}{48 U.S. at 37}).

\section{Procedural History}

After relocating to Massachusetts (establishing diversity of citizenship), Martin Luther filed an action of trespass \textit{quare clausum fregit} in October 1842 in the U.S. Circuit Court for the District of Rhode Island against Luther M. Borden and his fellow militiamen, all Rhode Island citizens (\href{https://supreme.justia.com/cases/federal/us/48/1/}{48 U.S. at 2, 34}). The defendants pleaded four pleas in justification, asserting that an armed insurrection existed, that the charter legislature had validly declared martial law, and that they entered Luther's house under military orders using no more force than necessary (\href{https://supreme.justia.com/cases/federal/us/48/1/}{48 U.S. at 2}).

At the November 1843 trial before Circuit Justice Joseph Story, Luther offered extensive documentary and testimonial evidence---including ballots, convention records, and the 1840 federal census---to prove that the People's Constitution had lawfully displaced the charter government prior to June 1842, rendering the defendants' martial-law defense void (\href{https://supreme.justia.com/cases/federal/us/48/1/}{48 U.S. at 11--18, 38}). Acting \textit{pro forma} by agreement of the parties to bring the constitutional questions to the Supreme Court, Justice Story excluded Luther's evidence, admitted the defendants' evidence, and instructed the jury that the charter government and its laws remained in full force and justified the defendants' acts (\href{https://supreme.justia.com/cases/federal/us/48/1/}{48 U.S. at 18--19, 38}). The jury returned a verdict for the defendants, and Luther sought review in the U.S. Supreme Court by writ of error (\href{https://supreme.justia.com/cases/federal/us/48/1/}{48 U.S. at 2, 18--19}). A companion suit brought by Martin's mother, \textit{Rachel Luther v. Luther M. Borden et al.}, came up on a \textit{pro forma} certificate of division presenting twenty-nine points (\href{https://supreme.justia.com/cases/federal/us/48/1/}{48 U.S. at 2, 47--48}). The cases were argued on January 21 and 24--28, 1848, and decided on January 3, 1849 (\href{https://supreme.justia.com/cases/federal/us/48/1/}{Justia}; \href{https://www.oyez.org/cases/1789-1850/48us1}{Oyez}).

\section{Judicial Votes}

The Supreme Court decided \textit{Luther v. Borden} by a \textbf{5--1} vote of the participating justices (\href{https://ballotpedia.org/Luther_v._Borden}{Ballotpedia}; \href{https://supreme.justia.com/cases/federal/us/48/1/}{48 U.S. at 48}).

\begin{itemize}
\item \textbf{Majority (5):} Chief Justice Roger B. Taney (author of the Opinion of the Court), joined by Justices John McLean, James Moore Wayne, Samuel Nelson, and Robert Cooper Grier (\href{https://ballotpedia.org/Luther_v._Borden}{Ballotpedia}).
\item \textbf{Dissenting in part / Concurring in part (1):} Justice Levi Woodbury filed an opinion styled as a dissent, expressly concurring with the majority that determining which state government was legitimate was a nonjusticiable political question, while dissenting solely from the majority's holding that the Rhode Island legislature could constitutionally declare statewide martial law (\href{https://supreme.justia.com/cases/federal/us/48/1/}{48 U.S. at 48, 51, 58--59}).
\item \textbf{Not Participating (3):} Justices John Catron, Peter Vivian Daniel, and John McKinley were absent due to ill health when the case was argued and did not participate (\href{https://supreme.justia.com/cases/federal/us/48/1/}{48 U.S. at 48 n.*}).
\end{itemize}

\section{Holding}

\textbf{Judgment:} Affirmed the Circuit Court's judgment for the defendants in \textit{Martin Luther v. Borden}. In the companion case, \textit{Rachel Luther v. Borden}, the Court dismissed the certificate of division for lack of jurisdiction and remanded to the Circuit Court because a trial court may not certify an entire case \textit{pro forma} (\href{https://supreme.justia.com/cases/federal/us/48/1/}{48 U.S. at 47--48}).

\textbf{Rule Announced:} Federal courts lack judicial authority to adjudicate which of two rival regimes is the lawful government of a state or whether a state government satisfies the Guarantee Clause of Article IV, Section 4; that determination is a political question committed to Congress (in seating state delegations and enforcing the Guarantee Clause) and to the President (in deciding whether to intervene against domestic violence under the Calling Forth Act of February 28, 1795), whose recognition binds the federal judiciary (\href{https://supreme.justia.com/cases/federal/us/48/1/}{48 U.S. at 42--45}; \href{https://www.law.cornell.edu/constitution-conan/article-3/section-2/clause-1/luther-v-borden-and-the-guarantee-clause}{Constitution Annotated}). Furthermore, a state may invoke martial law and use military force to suppress an armed insurrection too strong for civil authority, and officers acting under such orders may arrest suspected insurgents and enter homes on reasonable grounds, provided they use no more force than necessary (\href{https://supreme.justia.com/cases/federal/us/48/1/}{48 U.S. at 45--46}).

\textbf{Holding vs. Dicta:} The core holding encompasses three propositions: (1) federal courts must defer to Rhode Island's own courts and to the federal political branches regarding which regime constituted Rhode Island's lawful government in June 1842; (2) the President's recognition of the charter government under the Act of 1795 was conclusive on the courts; and (3) Rhode Island's temporary declaration of martial law justified Borden's warrantless entry using reasonable force (\href{https://supreme.justia.com/cases/federal/us/48/1/}{48 U.S. at 40, 43--46}). By contrast, Chief Justice Taney's broader suggestion that a permanent military government in a state would violate the Guarantee Clause and impose a duty on Congress to overthrow it is dicta, as Rhode Island's martial law was temporary (\href{https://supreme.justia.com/cases/federal/us/48/1/}{48 U.S. at 45}). Likewise, scholars such as Tara Leigh Grove have noted that Taney's sweeping statements rendering all Guarantee Clause questions nonjusticiable arguably went beyond the plaintiff's pleadings, which centered on popular sovereignty rather than a standalone Guarantee Clause claim (\href{https://www.law.cornell.edu/constitution-conan/article-3/section-2/clause-1/luther-v-borden-and-the-guarantee-clause}{Constitution Annotated}; \href{https://scholarship.law.wm.edu/cgi/viewcontent.cgi?article=2834&context=facpubs}{Grove, 90 N.Y.U. L. Rev. 1908, 1927--29 (2015)}).

\section{Analysis of Opinions}

\subsection{Majority Opinion (Chief Justice Taney)}

Chief Justice Taney rested the Court's refusal to adjudicate the legitimacy of the rival Rhode Island governments on four converging pillars: practical chaos, federalism and state-court deference, lack of judicial standards, and constitutional commitment to the political branches (\href{https://supreme.justia.com/cases/federal/us/48/1/}{48 U.S. at 38--45}).

First, Taney warned of catastrophic consequences if a federal court retroactively declared the charter government void during 1842--1843: every statute passed, tax collected, salary paid, and civil or criminal judgment rendered during that year would become a nullity, exposing state officers to civil and criminal liability (\href{https://supreme.justia.com/cases/federal/us/48/1/}{48 U.S. at 38--39}). Second, he emphasized that Rhode Island's own courts---including the Supreme Court of Rhode Island sitting under the undisputed 1843 Constitution in Thomas Dorr's treason trial---had held the charter government to be the lawful government in 1842; federal courts are bound to follow state tribunals on questions of state constitutional law (\href{https://supreme.justia.com/cases/federal/us/48/1/}{48 U.S. at 39--40}). Indeed, a state court could never declare the very government that created it illegitimate without simultaneously destroying its own judicial authority (\href{https://supreme.justia.com/cases/federal/us/48/1/}{48 U.S. at 40}).

Third, Taney demonstrated that the dispute lacked judicially manageable standards and procedures. A federal court has no authority to prescribe voter qualifications, conduct a census of freeholders, or adjudicate the eligibility of thousands of voters through witness testimony (\href{https://supreme.justia.com/cases/federal/us/48/1/}{48 U.S. at 41}). Worse, because Martin Luther's trespass suit was an action at common law governed by the Seventh Amendment, different juries in different cases could reach conflicting verdicts on whether the People's Constitution had been ratified, leaving state sovereignty perpetually unsettled (\href{https://supreme.justia.com/cases/federal/us/48/1/}{48 U.S. at 41--42}).

Fourth, Taney turned to Article IV, Section 4 (the Guarantee Clause), concluding that the Constitution committed both the guarantee of a republican form of government and protection against domestic violence exclusively to the political branches (\href{https://supreme.justia.com/cases/federal/us/48/1/}{48 U.S. at 42--43}). Congress exercises this authority when it decides whether to seat a state's senators and representatives, and Congress validly delegated to the President under the Act of February 28, 1795, the power to determine which faction represents the lawful state government when responding to an application for federal aid (\href{https://supreme.justia.com/cases/federal/us/48/1/}{48 U.S. at 42--44}). Citing \textit{Martin v. Mott}, 25 U.S. (12 Wheat.) 19 (1827), Taney held that President Tyler's recognition of the charter government was conclusive on the courts, just as executive recognition of foreign governments binds the judiciary (\href{https://supreme.justia.com/cases/federal/us/48/1/}{48 U.S. at 44--45}).

Finally, on the merits of Borden's trespass defense, Taney upheld the charter legislature's temporary declaration of martial law as a necessary measure of self-defense against armed insurrection (\href{https://supreme.justia.com/cases/federal/us/48/1/}{48 U.S. at 45--46}). Because it was a state of war, militia officers could arrest suspected insurgents and forcibly enter houses on reasonable grounds, though they would remain answerable for oppressive or willful excess force (\href{https://supreme.justia.com/cases/federal/us/48/1/}{48 U.S. at 45--46}).

\subsection{Concurring and Dissenting Opinion (Justice Woodbury)}

Justice Woodbury expressly agreed with the majority that the Court lacked jurisdiction to determine whether the People's Constitution had replaced the 1663 Charter (\href{https://supreme.justia.com/cases/federal/us/48/1/}{48 U.S. at 51}). He offered an additional separation-of-powers rationale: courts operate only \textit{beneath} established constitutions and laws to adjudicate private rights, never \textit{above} the people to arbitrate the formation of constitutions (\href{https://supreme.justia.com/cases/federal/us/48/1/}{48 U.S. at 51--53}). Entrusting life-tenured, unelected judges with final authority over political sovereignty would create an irresponsible oligarchy and strip the people of their birthright (\href{https://supreme.justia.com/cases/federal/us/48/1/}{48 U.S. at 52--54}).

Woodbury parted sharply from the majority, however, on the validity of Rhode Island's martial-law statute---a question he deemed squarely judicial because it involved private personal rights against warrantless home invasion (\href{https://supreme.justia.com/cases/federal/us/48/1/}{48 U.S. at 58--59}). Drawing on English constitutional history (Magna Carta, the Petition of Right of 1628, and the Bill of Rights of 1689) and American constitutional principles, Woodbury distinguished \textit{military law} (which governs only soldiers) from \textit{martial law} (which suspends civil law over the entire civilian population) (\href{https://supreme.justia.com/cases/federal/us/48/1/}{48 U.S. at 60--65}). He argued that neither the 1663 Charter nor any American state legislature possessed authority to suspend all civil liberties and subject civilians to warrantless military arrests while civil courts remained open and functioning (\href{https://supreme.justia.com/cases/federal/us/48/1/}{48 U.S. at 66--70}). In his view, officials who override ordinary law in an emergency should act at their peril and seek legislative indemnity afterward, rather than receiving advance judicial sanction for ``the reign of the bayonet'' (\href{https://supreme.justia.com/cases/federal/us/48/1/}{48 U.S. at 69}).

\section{Examples: Future Applications}

\begin{enumerate}
\item \textbf{Same-Side Application (Nonjusticiable Guarantee Clause Challenge to Direct Democracy):} A utility company challenges a state constitutional amendment adopted via ballot initiative, arguing that direct lawmaking by voters replaces representative democracy and therefore violates the Guarantee Clause. \textit{Why it falls on the same side:} Just as in \textit{Luther}, the Supreme Court held in \textit{Pacific States Telephone \& Telegraph Co. v. Oregon}, 223 U.S. 118 (1912), that whether a state's institutional structure remains ``republican'' under Article IV, Section 4, is a political question solely for Congress, not the federal courts (\href{https://www.law.cornell.edu/constitution-conan/article-3/section-2/clause-1/luther-v-borden-and-the-guarantee-clause}{Constitution Annotated}).

\item \textbf{Same-Side Application (Nonjusticiable Partisan Gerrymandering):} Voters sue in federal court alleging that a state's congressional map is so heavily gerrymandered for partisan advantage that it violates the Guarantee Clause and Article I by entrenching minority rule. \textit{Why it falls on the same side:} Echoing \textit{Luther}'s concern over a lack of judicially discoverable and manageable standards, the Supreme Court held in \textit{Rucho v. Common Cause}, 588 U.S. 684 (2019), that partisan gerrymandering claims present nonjusticiable political questions beyond federal judicial competence (\href{https://supreme.justia.com/cases/federal/us/588/18-422/}{Justia}).

\item \textbf{Opposite-Side Application (Justiciable Equal Protection Malapportionment Challenge):} Urban voters sue state election officials under the Fourteenth Amendment's Equal Protection Clause because a state legislature has refused to reapportion legislative districts for sixty years, leaving rural votes worth twenty times urban votes. \textit{Why it falls on the opposite side:} In \textit{Baker v. Carr}, 369 U.S. 186 (1962), and \textit{Reynolds v. Sims}, 377 U.S. 533 (1964), the Supreme Court distinguished \textit{Luther}, holding that while Guarantee Clause claims remain nonjusticiable, Equal Protection challenges to malapportionment are justiciable because the ``one person, one vote'' standard is judicially manageable and does not require questioning the state government's overall legitimacy (\href{https://www.law.cornell.edu/constitution-conan/article-3/section-2/clause-1/luther-v-borden-and-the-guarantee-clause}{Constitution Annotated}; \href{https://supreme.justia.com/cases/federal/us/377/533/}{Justia}).

\item \textbf{Opposite-Side Application (Justiciable Challenge to Excessive Force or Military Trials of Civilians):} During civil unrest, state militia officers acting under a governor's emergency proclamation arrest a peaceful journalist with no connection to the unrest, destroy her printing press out of personal animus, and try her before a military tribunal while state courts remain open. \textit{Why it falls on the opposite side:} Even under Chief Justice Taney's majority opinion in \textit{Luther}, officers who exercise emergency power ``for the purposes of oppression'' or inflict willful injury remain answerable in court (\href{https://supreme.justia.com/cases/federal/us/48/1/}{48 U.S. at 46}), and subsequent precedent vindicating Justice Woodbury's dissent holds that military trials of civilians when civil courts are open are judicially reviewable and unconstitutional.

\item \textbf{Fence-Sitter Application (Federal Commandeering and the Guarantee Clause):} Congress enacts a statute compelling state legislatures to pass specific regulatory laws or take title to hazardous waste, and a state sues under both the Tenth Amendment and the Guarantee Clause, arguing that federal coercion destroys the state's republican accountability to its voters. \textit{Why it is a fence-sitter:} In \textit{New York v. United States}, 505 U.S. 144, 184--85 (1992), Justice O'Connor explicitly questioned whether \textit{Luther}'s blanket nonjusticiability rule was overbroad and suggested that ``perhaps not all claims under the Guarantee Clause present nonjusticiable political questions,'' while avoiding a definitive ruling because the statute failed on Tenth Amendment grounds (\href{https://www.law.cornell.edu/constitution-conan/article-3/section-2/clause-1/luther-v-borden-and-the-guarantee-clause}{Constitution Annotated}). Whether a federal-interference Guarantee Clause claim is justiciable remains unresolved.
\end{enumerate}

\section{Critique}

\textit{Luther v. Borden} occupies a pivotal and fiercely contested place in constitutional scholarship as the foundational precedent of the modern political question doctrine (\href{https://www.law.cornell.edu/constitution-conan/article-3/section-2/clause-1/luther-v-borden-and-the-guarantee-clause}{Constitution Annotated}).

From an \textbf{originalist and textualist} perspective, scholars divide sharply over Taney's reading of Article IV, Section 4. On one hand, scholars like Ryan C. Williams and Heritage Foundation commentators note that the Clause's text assigns the duty to ``[t]he United States''---not specifically to Congress---leaving open whether courts share in enforcing the guarantee, and historical research shows the Framers primarily aimed to prevent monarchical or aristocratic usurpations while preserving state flexibility (\href{https://harvardlawreview.org/wp-content/uploads/2018/12/602-688_Online.pdf}{Williams, 132 Harv. L. Rev. 602 (2018)}; \href{https://constitution.heritage.org/essays/a4-s4/}{Heritage Guide}). Tara Leigh Grove further argues that nineteenth-century ``political question'' cases like \textit{Luther} originally involved deference to specific factual determinations made by the political branches (such as recognizing a government), rather than a sweeping jurisdictional bar stripping Article III courts of authority to interpret constitutional text (\href{https://scholarship.law.wm.edu/cgi/viewcontent.cgi?article=2834&context=facpubs}{Grove, 90 N.Y.U. L. Rev. 1908 (2015)}). Conversely, other formalists and textualists defend \textit{Luther} and call for reinvigorating the political question doctrine---especially where constitutional text commits structural or electoral matters to political branches---to restrain unelected judges from making unguided policy choices (\href{https://scholarlycommons.law.northwestern.edu/cgi/viewcontent.cgi?article=1477&context=nulr}{Northwestern Univ. L. Rev.}; \href{https://www.law.georgetown.edu/public-policy-journal/wp-content/uploads/sites/23/2022/09/GT-GLPP220051-1.pdf}{Georgetown J.L. \& Pub. Pol'y}).

From a \textbf{progressive and doctrinal} perspective, prominent scholars such as Erwin Chemerinsky criticize \textit{Luther} and its progeny for effectively erasing the Guarantee Clause from the Constitution (\href{https://scholarship.law.duke.edu/cgi/viewcontent.cgi?article=1696&context=faculty_scholarship}{Chemerinsky, Duke Law Scholarship}). Chemerinsky contends that refusing to enforce the Guarantee Clause abdicates the judiciary's \textit{Marbury} duty to say what the law is, leaves democratic deficits (such as severe disenfranchisement or anti-majoritarian entrenchment) without judicial remedy when Congress refuses to act, and forces courts to stretch other clauses like Equal Protection to do the work the Guarantee Clause was written to perform (\href{https://scholarship.law.duke.edu/cgi/viewcontent.cgi?article=1696&context=faculty_scholarship}{Chemerinsky}; \href{https://arizonalawreview.org/pdf/62-1/62arizlrev183.pdf}{62 Ariz. L. Rev. 183 (2020)}). Historians have also pointed out the antebellum subtext of Taney's opinion: declaring the Guarantee Clause nonjusticiable conveniently insulated Southern slave states---whose restricted suffrage and chattel slavery abolitionists attacked as un-republican---from federal judicial scrutiny (\href{https://www.law.cornell.edu/constitution-conan/article-3/section-2/clause-1/luther-v-borden-and-the-guarantee-clause}{Constitution Annotated}; \href{https://law.utexas.edu/faculty/uploads/publication_files/luther-v.-borden--a-taney-court-mystery-solved.pdf}{Texas Law}). Meanwhile, Justice Woodbury's dissent on martial law is widely praised across ideological lines as a prescient defense of civil liberties against unchecked emergency powers.

At bottom, \textit{Luther} reflects a deep tension among \textbf{competing constitutional values}: judicial fidelity to constitutional text and individual remedies versus institutional competence, democratic self-governance, and stability; and state order during domestic insurrection versus the rule of law and protection against military despotism.

\section{Key Quotations}

\begin{quote}
``Undoubtedly the courts of the United States have certain powers under the Constitution and laws of the United States which do not belong to the State courts. But the power of determining that a State government has been lawfully established, which the courts of the State disown and repudiate, is not one of them.''
\end{quote}
\noindent --- Chief Justice Roger B. Taney, \textit{Luther v. Borden}, 48 U.S. (7 How.) 1, 40 (1849) (\href{https://supreme.justia.com/cases/federal/us/48/1/}{Justia}; \href{https://www.courtlistener.com/opinion/86466/luther-v-borden/}{CourtListener}).

\begin{quote}
``Under this article of the Constitution, it rests with Congress to decide what government is the established one in a State. For as the United States guarantee to each State a republican government, Congress must necessarily decide what government is established in the State before it can determine whether it is republican or not.''
\end{quote}
\noindent --- Chief Justice Roger B. Taney, \textit{Luther v. Borden}, 48 U.S. (7 How.) 1, 42 (1849) (\href{https://supreme.justia.com/cases/federal/us/48/1/}{Justia}; \href{https://www.courtlistener.com/opinion/86466/luther-v-borden/}{CourtListener}).

\begin{quote}
``If the judicial power extends so far, the guarantee contained in the Constitution of the United States is a guarantee of anarchy, and not of order.''
\end{quote}
\noindent --- Chief Justice Roger B. Taney, \textit{Luther v. Borden}, 48 U.S. (7 How.) 1, 43 (1849) (\href{https://supreme.justia.com/cases/federal/us/48/1/}{Justia}; \href{https://www.courtlistener.com/opinion/86466/luther-v-borden/}{CourtListener}).

\begin{quote}
``And unquestionably a State may use its military power to put down an armed insurrection too strong to be controlled by the civil authority. The power is essential to the existence of every government, essential to the preservation of order and free institutions, and is as necessary to the States of this Union as to any other government.''
\end{quote}
\noindent --- Chief Justice Roger B. Taney, \textit{Luther v. Borden}, 48 U.S. (7 How.) 1, 45 (1849) (\href{https://supreme.justia.com/cases/federal/us/48/1/}{Justia}; \href{https://www.courtlistener.com/opinion/86466/luther-v-borden/}{CourtListener}).

\begin{quote}
``Our power begins after theirs ends. Constitutions and laws precede the judiciary, and we act only under and after them, and as to disputed rights beneath them, rather than disputed points in making them.''
\end{quote}
\noindent --- Justice Levi Woodbury, dissenting in part, \textit{Luther v. Borden}, 48 U.S. (7 How.) 1, 52 (1849) (\href{https://supreme.justia.com/cases/federal/us/48/1/}{Justia}; \href{https://www.courtlistener.com/opinion/86466/luther-v-borden/}{CourtListener}).

\end{document}
